\section{Notation}\label{sec:notation}

Table~\ref{tab:notation} collects every symbol of the report with its unit. Four conventions hold throughout and are worth
stating before the methods, because the engine depends on them and a reimplementation that changes one of them gives
different numbers.

\emph{Size classes.} The sieve grid runs from coarse to fine. Edge $e_k$ is $256/(\sqrt2)^{k}$~mm for $k=0,\dots,28$, so
there are $K=28$ classes and the finest edge is $256/2^{14}=0.015625$~mm. Class $k$ lies between the upper edge $e_k$ and
the lower edge $e_{k+1}$ and has the geometric mid-size $\bar d_k=\sqrt{e_k e_{k+1}}$. Because the classes run coarse to
fine, a matrix whose entry $(i,j)$ moves mass from class $j$ to class $i$ is strictly lower triangular when mass only moves
to finer classes, and the breakage matrix of Section~\ref{sec:m-breakage} has this form. Vectors indexed by class are bold
($\mathbf f$, $\mathbf p$, $\mathbf c$); a diagonal matrix built from a vector is the upright capital ($\mathrm C=\mathrm{diag}(\mathbf c)$).

\emph{Gradations.} A gradation is stored either as mass per class (a vector that sums to one) or as cumulative fraction
passing at each edge (a vector that starts at one and decreases). The size at a passing fraction ($P_{80}$, $P_{50}$,
$P_{20}$, $F_{80}$) is read by linear interpolation in the logarithm of the size between the two edges that bracket the
fraction. Capital $F_{80}$ is the feed $80\%$ passing size as the engine reads it on its grid; lower-case $f_{80}$ is the
measured value of a survey. They differ when the feed is coarser than the $256$~mm top of the grid.

\emph{Units.} Sizes are in millimetres in the engine and in micrometres inside the two energy laws (Bond and Morrell), as
those laws are printed; throughput is in tonnes per hour; power in kilowatts; specific energy in kWh per tonne; speed in
revolutions per minute; angles in degrees. Percentages are written with the percent sign; a difference between two
percentage errors is written in points.

\emph{Survey indices and folds.} Survey $i$ runs from $1$ to $N_s=10$ in the order of the published table. A fold of the
leave-one-survey-out evaluation is indexed by its held-out survey; $\mathcal T_{-i}$ is the set of the other nine surveys,
and a quantity fitted on them carries the subscript $-i$. A standardised variable is written $z$, with the mean and
standard deviation of the training surveys of the fold. In CrusherBayes, $\boldsymbol\theta$ is the vector of the eight
engine scalars and the two noise scales, and $\mathbf d_t$ the three discrepancy coefficients of target $t$.

\begin{table}[!t]
\centering
\caption{Notation. Every symbol used in the equations of the report, its meaning and its unit; a dash marks a
dimensionless quantity. Symbols local to one derivation are defined where they appear.}
\label{tab:notation}
\footnotesize
\begin{tabular}{@{}l p{0.62\textwidth} l@{}}
\toprule
Symbol & Meaning & Unit \\
\midrule
$K$, $k$ & number of size classes ($28$) and class index, coarse to fine & -- \\
$e_k$, $\bar d_k$ & upper edge and geometric mid-size of class $k$ & mm \\
$\mathbf f$, $\mathbf p$ & feed and product mass per class (each sums to one) & -- \\
$R(x)$ & cumulative fraction passing size $x$ & -- \\
$x_{63}$, $n$ & Rosin-Rammler characteristic size and uniformity exponent & mm, -- \\
$F_{80}$, $f_{80}$ & feed $80\%$ passing size, on the engine grid and as measured & mm \\
$P_{80}$, $P_{50}$, $P_{20}$ & product sizes passing $80$, $50$ and $20\%$ & mm \\
$s$ (CSS), $T$, OSS & closed-side setting, eccentric throw, open-side setting $s+T$ & mm \\
$T_{ref}$ & design throw of a machine class & mm \\
$N$, $N_{opt}$ & eccentric speed and the speed of maximum capacity & rpm \\
$A$, $b$, $A\!\times\!b$ & JKMRC ore competence parameters and their product & \%, t/kWh, -- \\
$E_{cs}$ & specific comminution energy per nip & kWh/t \\
$t_{10}$ & fraction of progeny passing one tenth of the parent size (written in percent where stated) & -- \\
$B(u)$, $\varphi$, $\gamma$, $\beta$ & Austin cumulative appearance function of relative size $u$, its weight and slopes & -- \\
$\mathrm B$, $B_{ij}$ & breakage matrix, fraction of broken class-$j$ mass reporting to class $i$ & -- \\
$C(d)$, $\mathbf c$, $\mathrm C$ & classification probability of size $d$, its vector and diagonal matrix & -- \\
$K_1$, $K_2$, $K_3$ & classification sizes and exponent & mm, mm, -- \\
$\mathbf x$ & internal mass flow of the breakage zone per unit feed & -- \\
$\kappa$ & conditioning proxy of $\mathrm I-\mathrm B\mathrm C$ & -- \\
$Q$, $Q_{ref}$, $Q_T$ & throughput, reference capacity, calibrated capacity base & t/h \\
$g_{ref}$ & reference discharge gap of a machine class & mm \\
$W_i$, $W$ & Bond work index and Bond specific energy & kWh/t \\
$P_d$, $P_c$, $P$ & Bond net draw, current-based calibrated power, power & kW \\
$\zeta$, $P_n$ & slope and no-load term of the current-based power form & --, kW \\
$M_{ic}$, $W_c$, $S_c$ & Morrell crushing work index, crusher specific energy, hardness correction & kWh/t, kWh/t, -- \\
$x_1$, $x_2$ & feed and product $80\%$ passing sizes in the Morrell equation & $\mu$m \\
$\alpha$, $\mu_f$ & nip angle and friction coefficient of the grip limit & deg, -- \\
$\mathcal D$, $N_s$ & set of surveys and their number ($10$) & -- \\
$\mathcal T_{-i}$ & training surveys of the fold that holds out survey $i$ & -- \\
$y_{t,i}$, $\hat y_{t,i}$ & measured value and held-out prediction of target $t$ at survey $i$ & t/h or kW \\
$\sigma$, $\sigma_{-i}$ & residual standard deviation, and that of the nine other surveys & t/h or kW \\
$z_{\mathrm{CSS}}$, $z_{f_{80}}$ & setting and feed size standardised on the training surveys & -- \\
$\boldsymbol\theta$ & CrusherBayes engine scalars and log noise scales & mixed \\
$\delta_t$, $\mathbf d_t$, $\tau_t$ & discrepancy of target $t$, its coefficients and their prior scale & t/h or kW \\
$\sigma_t$ & observation noise of target $t$ & t/h or kW \\
$\hat\tau$ & integrated autocorrelation time of a chain & steps \\
$\Delta s$, $r$, $t_{out}$ & setting opening to a target, drift rate, hours to the target & mm, mm/h, h \\
\bottomrule
\end{tabular}
\end{table}
